\documentclass[11pt]{article}
\usepackage[margin=1in]{geometry}
\usepackage{amsmath,amssymb}
\usepackage{graphicx}
\usepackage{booktabs}
\usepackage{hyperref}
\usepackage{fancyhdr}
\usepackage{caption}
\captionsetup{font=small}
\pagestyle{fancy}
\fancyhf{}
\lhead{Project Report}
\rhead{Numerical Study of Inflationary Cosmology}
\rfoot{\thepage}
\setlength{\headheight}{14pt}

\title{Numerical Study of Inflationary Cosmology:\\
\large Scalar and Tensor Perturbation Spectra Across Single-Field Potentials}
\author{Rishi Ramachandran}
\date{Project Report}

\begin{document}
\maketitle

\begin{abstract}
This project developed a numerical framework for evolving single-field inflationary
backgrounds and both scalar and tensor perturbations, in order to compute the primordial
curvature and gravitational-wave power spectra for a range of inflationary potentials.
Starting from the Mukhanov--Sasaki formalism, the framework evolves the background in
e-fold time, initializes perturbations deep inside the Hubble radius with Bunch--Davies
initial conditions, follows them into the super-Hubble regime, and extracts the frozen
scalar and tensor power spectra. The framework was validated on quadratic inflation
against slow-roll analytic expectations, then extended to seven further potentials
(quartic, hilltop, Starobinsky, natural, symmetry-breaking, a quadratic+quartic mixture,
and the T-model $\alpha$-attractor), all calibrated to a shared benchmark of
$N_{\rm end}=65$ e-folds and pivot amplitude $\mathcal{P}_\mathcal{R}(k_*)=2.05\times10^{-9}$.
Systematic parameter sweeps across each potential family were used to explore how the
predicted $(n_s,r)$ pair moves as a function of the potential's free parameters, revealing
a classification of parameters into \emph{scale-type} (rescaling $\phi$ at fixed potential
shape) and \emph{shape-type} (changing the functional form itself). This classification,
and the resulting exponential- versus linear-track behaviour it produces, was observed and
numerically verified across seven independent sweeps, and is flagged here as a pattern
worth further analytic and observational investigation.
\end{abstract}

\section{Introduction}
Inflation explains the horizon and flatness problems through a period of accelerated
expansion, and predicts a nearly scale-invariant spectrum of scalar curvature perturbations
together with a spectrum of tensor perturbations (primordial gravitational waves). Both are
computed by evolving the homogeneous inflaton background and the linear perturbations
around it, using the Mukhanov--Sasaki equation.

This project began as a single-potential numerical implementation of that pipeline
(quadratic inflation, validated against slow-roll expectations) and has since been
extended substantially: to seven additional potentials, to the tensor sector and the
tensor-to-scalar ratio $r$, and to systematic parameter sweeps within each potential
family. The unifying design principle throughout is that the potential $V(\phi)$ is
supplied to a fixed background/perturbation solver, so that extending to a new model is a
local change (new $V,V_{,\phi}$) rather than a rewrite of the numerics. All theoretical
background used here follows Baumann's \emph{Cosmology} lecture notes~\cite{baumann},
equation numbers below are given for the
project's own numbering, with the corresponding Baumann equation noted where directly
quoted.

\section{Theoretical Framework}

\subsection{Background dynamics}
For a homogeneous inflaton in a flat FRW background, working in e-fold time
$N\equiv\ln a$ and reduced Planck units, the Friedmann constraint and Klein--Gordon
equation become
\begin{equation}
H^2=\frac{V(\phi)}{3-\tfrac12\phi_N^2},\qquad
\phi_{NN}=-\Big(3-\tfrac12\phi_N^2\Big)\phi_N-\Big(6-\phi_N^2\Big)\frac{V_{,\phi}}{2V},
\label{eq:bg}
\end{equation}
with $\phi_N\equiv d\phi/dN$ [Baumann eqs.\ (2.3.24), (2.3.27), recast in $N$-time].
The first Hubble slow-roll parameter is $\varepsilon_1=\tfrac12\phi_N^2$, and inflation
ends when $\varepsilon_1=1$ ($\phi_N=\sqrt2$); this is the numerical stopping event used
throughout. The potential slow-roll parameters, evaluated directly from $V(\phi)$, are
\begin{equation}
\varepsilon_v(\phi)=\tfrac12\big(V_{,\phi}/V\big)^2,\qquad
\eta_v(\phi)=V_{,\phi\phi}/V,
\label{eq:sr}
\end{equation}
[Baumann eq.\ (2.3.37)], used below as the analytic slow-roll cross-check on every
numerical result.

\subsection{Scalar and tensor perturbations}
The gauge-invariant comoving curvature perturbation $\mathcal{R}_k$ obeys the
Mukhanov--Sasaki equation, which in $e$-fold time and with $z\equiv a\phi_N$ reads
\begin{equation}
\mathcal{R}_{k,NN}+\Big[1-\tfrac12\phi_N^2+2\frac{z_N}{z}\Big]\mathcal{R}_{k,N}
+\Big(\frac{k}{aH}\Big)^2\mathcal{R}_k=0,
\label{eq:ms-scalar}
\end{equation}
[Baumann eq.\ (6.1.9), $z$-form]. Tensor perturbations $h_k$ (per polarization) obey the
same equation with $z\to a$ and no $\phi_N$-dependent friction term beyond the standard
$3H$ Hubble drag. Both are complex, split into real/imaginary parts, and initialized deep
inside the horizon ($k/aH\simeq100$) with Bunch--Davies vacuum initial conditions
[Baumann eq.\ (6.3.41)], then integrated to $k/aH\simeq10^{-5}$, well outside the horizon,
where both $\mathcal{R}_k$ and $h_k$ freeze to constants. The dimensionless power spectra
are then
\begin{equation}
\mathcal{P}_\mathcal{R}(k)=\frac{k^3}{2\pi^2}|\mathcal{R}_k|^2,\qquad
\mathcal{P}_h(k)=8\,\frac{k^3}{2\pi^2}|h_k|^2,\qquad
r(k)\equiv\frac{\mathcal{P}_h(k)}{\mathcal{P}_\mathcal{R}(k)},
\label{eq:spectra}
\end{equation}
with the tensor prefactor of 8 (two polarizations, and normalization of $h_k$ to the
physical metric perturbation) calibrated against the slow-roll result $r=16\varepsilon_v$
[Baumann eq.\ (6.5.65)]. The scalar spectral index is extracted numerically as the local
log-slope $n_s-1\equiv d\ln\mathcal{P}_\mathcal{R}/d\ln k$ near the pivot, and compared to
the analytic slow-roll formula $n_s-1=-6\varepsilon_v+2\eta_v$ [Baumann eq.\ (6.4.57)]
evaluated at the field value $N_*=50$ e-folds before the end of inflation. The tensor
spectral index $n_t\equiv d\ln\mathcal{P}_h/d\ln k$ is extracted the same way, and used
below to test the standard single-field consistency relation $r=-8n_t$.

\section{Numerical Framework and Shared Benchmark}
Every potential in this project is evolved through the same pipeline: (1) specify
$V,V_{,\phi}$; (2) root-find the initial field value $\phi_i$ on the \emph{exact}
background ODE~\eqref{eq:bg} (not the slow-roll estimate) so that inflation ends at
exactly $N_{\rm end}=65$; (3) tune the potential's overall scale so the pivot-scale
amplitude matches $\mathcal{P}_\mathcal{R}(k_*)=2.05\times10^{-9}$ at $k_*=0.05\,{\rm
Mpc^{-1}}$; (4) locate the sub-Hubble ($k/aH=100$, i.e.\ \texttt{Nic\_cond=100}) and
super-Hubble ($k/aH=10^{-5}$) boundaries for each mode; (5) integrate
Eq.~\eqref{eq:ms-scalar} (and its tensor counterpart) with absolute tolerance
$\texttt{atol}=10^{-13}$; (6) evaluate the frozen spectra via Eq.~\eqref{eq:spectra}.
$N_*=50$ (the pivot mode's horizon-crossing e-fold, counted \emph{before the end of
inflation}) is fixed by construction across every run, while $N_{\rm end}=65$ is a
tunable target common to all potentials -- this distinction matters because, as shown
below, the physics at the pivot depends only on $N_*$, not on $N_{\rm end}$, for
potentials with a single slow-roll attractor trajectory.

\subsection{Validation: quadratic inflation}
As in the original single-potential version of this project, $V=\tfrac12 m^2\phi^2$
remains the principal validation case, since both its background and its spectral
predictions have closed-form slow-roll expectations
($\varepsilon_v=\eta_v=2/\phi^2$, $N(\phi)=\phi^2/4-1/2$). Re-targeted to the shared
$N_{\rm end}=65$ benchmark, root-finding gives $\phi_i=16.055\,M_{pl}$ and
$m=6.973\times10^{-6}\,M_{pl}$. Figure~\ref{fig:validation} shows the resulting pivot-mode
curvature perturbation -- oscillating while sub-Hubble, freezing after horizon exit -- and
the resulting scalar spectrum compared with the slow-roll power-law benchmark.
Table~\ref{tab:main} (quadratic row) shows the numerical $(n_s,r)$ agreeing with the
analytic slow-roll values to within $\lesssim0.5\%$ ($n_s$) and $\lesssim3\%$ ($r$),
consistent with the expected size of slow-roll corrections at $N_*=50$.

\begin{figure}[htbp]
\centering
\includegraphics[width=\textwidth]{figs/fig_validation.pdf}
\caption{Validation case (quadratic inflation, $N_{\rm end}=65$). Left: real and
imaginary parts of the pivot-mode curvature perturbation, oscillating sub-Hubble and
freezing on super-Hubble scales. Right: numerical scalar power spectrum compared with the
slow-roll power-law form.}
\label{fig:validation}
\end{figure}

\section{Extension to Seven Further Potentials}
The same pipeline, unmodified, was applied to seven additional potentials: quartic
$V=\tfrac14\lambda\phi^4$; hilltop $V=V_0[1-(\phi/\mu)^2]$; Starobinsky
$V=V_0(1-e^{-\sqrt{2/3}\,\phi})^2$; natural inflation $V=\Lambda^4(1-\cos(\phi/f))$;
symmetry breaking $V=V_0(1-(\phi/\mu)^2)^2$; a quadratic+quartic mixture
$V=\tfrac12\phi^2+\tfrac{g}{4}\phi^4$ (rescaled to the pivot amplitude); and the T-model
$\alpha$-attractor $V=V_0\tanh^{2n}(\phi/\sqrt{6\alpha})$. All eight are shown together
in Figure~\ref{fig:background} (background evolution and end-of-inflation sharpness) and
Figure~\ref{fig:spectra} (scalar and tensor spectra), all calibrated to the shared
benchmark. Table~\ref{tab:main} collects the numerical results for one representative
parameter choice per potential.

\begin{figure}[htbp]
\centering
\includegraphics[width=\textwidth]{figs/fig_background.pdf}
\caption{Background field evolution (left) and the first Hubble slow-roll parameter
aligned at the end of inflation (right), for all eight potentials at the shared
$N_{\rm end}=65$ benchmark. Nonlinear potentials (hilltop, symmetry breaking,
Starobinsky) show a visibly sharper $\varepsilon_1(N)\to1$ transition than the
monomial/linear potentials.}
\label{fig:background}
\end{figure}

\begin{figure}[htbp]
\centering
\includegraphics[width=\textwidth]{figs/fig_spectra.pdf}
\caption{Scalar (left) and tensor (right) power spectra for all eight potentials.}
\label{fig:spectra}
\end{figure}

\begin{table}[htbp]
\centering
\small
\begin{tabular}{lrrrrrr}
\toprule
Potential & $\phi_i\,[M_{pl}]$ & $n_s$ & $r$ & $n_t$ & $-8n_t$ & diff.\ from $r$ \\
\midrule
Quadratic            & 16.055 & 0.9597 & 0.15823 & $-0.02040$ & 0.16321 & $-3.1\%$ \\
Quartic               & 22.796 & 0.9393 & 0.31321 & $-0.04079$ & 0.32632 & $-4.0\%$ \\
Hilltop ($\mu=5.5$)   & 0.0513 & 0.8703 & 0.00056 & $-0.00008$ & 0.00064 & $-12.9\%$ \\
Starobinsky           & 5.518  & 0.9615 & 0.00426 & $-0.00055$ & 0.00443 & $-3.9\%$ \\
Natural ($f=5$)       & 12.901 & 0.9476 & 0.04926 & $-0.00648$ & 0.05182 & $-4.9\%$ \\
Symm.\ break.\ ($\mu=7$)& 0.0246& 0.8409 & 0.00030 & $-0.00004$ & 0.00035 & $-15.5\%$ \\
Mix ($g=0.1$)         & 21.862 & 0.9367 & 0.32886 & $-0.04289$ & 0.34315 & $-4.2\%$ \\
T-model ($\alpha=5$)  & 9.825  & 0.9602 & 0.02070 & $-0.00269$ & 0.02154 & $-3.9\%$ \\
\bottomrule
\end{tabular}
\caption{Numerical results for all eight potentials at $N_{\rm end}=65$,
$\mathcal{P}_\mathcal{R}(k_*)=2.05\times10^{-9}$. $n_s,n_t$ are local log-slope fits of
the numerical spectra near $k_*$; $r$ is $\mathcal{P}_h/\mathcal{P}_\mathcal{R}$ at the
pivot. The last column tests the single-field consistency relation $r=-8n_t$: all eight
potentials violate it at the several-percent level expected from next-to-leading-order
slow-roll corrections, with hilltop and symmetry breaking -- the two potentials with the
largest $|\eta_v/\varepsilon_v|$ at the pivot -- showing the largest violation.}
\label{tab:main}
\end{table}

A consistent pattern across the two ``hilltop-like'' potentials (hilltop and symmetry
breaking) is that they show the sharpest $\varepsilon_1(N)$ transition near end of
inflation (Fig.~\ref{fig:background}, right panel) and the largest deviation from the
consistency relation at the pivot (Table~\ref{tab:main}). Both are properties of the
same underlying feature: near $\phi=0$ these potentials are locally quadratic in $\phi^2$
around a maximum, which makes $|\eta_v|$ large relative to $\varepsilon_v$ at horizon
crossing, and next-to-leading-order slow-roll corrections (not captured by the
leading-order $r=-8n_t$ relation) scale with exactly this ratio.

\section{Parameter Sweeps and a Preliminary Classification}
\label{sec:classification}
For each potential family, a sweep over its free parameter(s) was carried out: for each
value, $\phi_i$ is re-root-found to hold $N_{\rm end}=65$ fixed and the amplitude is
recalibrated, isolating the effect of the parameter on $(n_s,r)$ alone. Six sweeps were
performed: hilltop $\mu$, natural $f$, symmetry-breaking $\mu$ (all rescale $\phi$ at
fixed potential shape); monomial power $p$ and the quadratic+quartic coupling ratio $g$
(both change the functional form of $V$); and the T-model, swept separately over the
curvature scale $\alpha$ and the exponent $n$.

\textbf{Observed pattern.} Sweeping the scale-type parameters (hilltop $\mu$, natural
$f$, symmetry-breaking $\mu$, T-model $\alpha$) traces out a track in the $(n_s,r)$ plane
that is well fit by an exponential/log-linear relation between $r$ and $n_s$
($R^2=0.9999$ for symmetry breaking and T-model $\alpha$; see
Fig.~\ref{fig:classification}). Sweeping the monomial power $p$ instead traces an
\emph{exactly linear} relation, consistent with the closed-form slow-roll result derived
in the accompanying primer (Appendix, Cumulative Review Problem 2): for
$V=\mu^{4-p}\phi^p$, $n_s-1=-(2+p)/2N_*$ and $r=4p/N_*$ exactly eliminate to a line in
$(n_s,r)$ at fixed $N_*$. Sweeping the quadratic+quartic coupling $g$ instead, $r$ moves
\emph{non-monotonically} with $g$ (rising, then falling slightly for $g\gtrsim0.05$), and
the log-linear fit is markedly worse ($R^2=0.93$) -- this potential does not fall cleanly
into either the ``rescale $\phi$'' or ``exact monomial'' bucket, since $g$ changes the
\emph{relative weight} of two different functional terms rather than rescaling a single
one. Finally, sweeping the T-model exponent $n$ at fixed $\alpha$ leaves both $n_s$ and
$r$ almost unchanged ($n_s\in[0.9588,0.9596]$, $r\in[0.00475,0.00493]$ across
$n\in[0.5,7]$) -- consistent with the known $\alpha$-attractor property that predictions
are controlled almost entirely by $\alpha$, with only weak sensitivity to the
tanh power $n$.

\begin{figure}[htbp]
\centering
\includegraphics[width=\textwidth]{figs/fig_classification.pdf}
\caption{$(n_s,r)$ tracks traced out by sweeping one parameter at a time within each
potential family, all at fixed $N_{\rm end}=65$, shown as individual panels. Circle/triangle
markers: parameters that rescale $\phi$ at fixed potential shape (``scale-type'').
Square markers: parameters that change the functional form of $V$ (``shape-type'').
This classification was observed directly in the numerical sweep data across all seven
parameters shown.}
\label{fig:classification}
\end{figure}

\textbf{Classification.} Based on this pattern, a potential's free parameters can be split
into two kinds: (i) \emph{scale-type} parameters, which rescale $\phi$ while holding the
shape of $V(\phi/{\rm scale})$ fixed, and which produce smooth, exponential-looking tracks
in $(n_s,r)$; and (ii) \emph{shape-type} parameters, which change the functional form of
$V$ itself, and which produce either an exact linear relation (when the change is a clean
power-law deformation, as for the monomial power $p$) or non-monotonic, harder-to-fit
behaviour (when two functional terms are traded off against each other, as for the
mixing parameter $g$). The T-model $n$-sweep suggests a possible third category --
parameters to which the $(n_s,r)$ prediction is largely insensitive -- though with only
one example ($\alpha$-attractors are already known analytically to have this property)
this is not yet confirmed as a general pattern and is a natural next thing to check.
This classification was observed and numerically verified across all seven sweeps; the
monomial case additionally has an accompanying analytic derivation (Sec.~3.7--B of the
accompanying primer) that matches the numerical track exactly, which is encouraging for
the scale/shape split more broadly, though that broader argument has not yet been derived
analytically for the other potentials. Extending this to an analytic argument for why scale-type
parameters generically produce a specific track shape, and confirming or refuting the
apparent third category, are the natural next steps (Sec.~\ref{sec:limitations}).

\section{Discussion}
The central outcome of this phase of the project is a validated, potential-agnostic
numerical pipeline that computes both scalar and tensor spectra for any single-field
potential, together with a first systematic look at how those predictions move as a
potential's parameters vary. Two things are on comparatively solid ground: the scalar
sector reproduces the closed-form quadratic benchmark to sub-percent accuracy, and the
qualitative pattern that nonlinear potentials (hilltop-type) produce sharper end-of-inflation
transitions and larger consistency-relation violations is visible directly in the raw
$\varepsilon_1(N)$ and $r$ vs.\ $-8n_t$ data, not inferred. The classification of
parameters into scale- and shape-type (Sec.~\ref{sec:classification}) is a pattern that
was directly observed and numerically verified across seven independent sweeps, one of
which (monomial) also has an analytic underpinning that matches it exactly. It is the
leading candidate for the project's next phase, and the natural next step is to push it
from an observed, verified pattern to a general analytic argument.

\section{Limitations and Future Development}
\label{sec:limitations}
\begin{itemize}
\item \textbf{Classification framework needs an analytic argument.} The scale/shape split
was observed and numerically verified across seven sweeps, but beyond the monomial case
no analytic argument has yet been constructed for why scale-type parameters should
produce an exponential-looking $(n_s,r)$ track, or for the apparent third (insensitive)
category seen only in the T-model $n$-sweep. This is the most important open item, and
the most promising direction for the next phase of the project.
\item \textbf{Consistency-relation check is leading order only.} $r=-8n_t$ is a
leading-order slow-roll relation; the observed several-percent violations are consistent
with next-to-leading-order corrections, but this project has not computed those
corrections independently, so the size of the violation is reported without a
next-to-leading-order comparison to confirm it is not, instead, a numerical-tolerance
artifact.
\item \textbf{Arbitrary potentials are not yet systematically validated.} As before,
quadratic inflation remains the only potential with a fully closed-form analytic
benchmark; all others are cross-checked only against their own leading-order slow-roll
formulas evaluated at the numerical background trajectory.
\item \textbf{Symmetry-breaking sweep numerics.} The $\mu=5$ point in the symmetry-breaking
sweep triggered floating-point overflow warnings during root-finding (Sec.~5); the
tabulated result is still self-consistent (probed against $\phi_N$ boundary conditions)
but the reliability of the $\mu\to$ small end of that track has not been separately
stress-tested.
\item \textbf{Only one $N_{\rm end}$ benchmark tested.} All results here use
$N_{\rm end}=65$; since $N_*=50$ is fixed by construction, changing $N_{\rm end}$ affects
only the trajectory between horizon crossing and the end of inflation, which is why the
quadratic pivot values matched exactly between an earlier $N_{\rm end}=60$ pass and the
current one -- but this has only been explicitly checked for quadratic, not for the other
seven potentials.
\end{itemize}

\section{Research Experience}
As with the earlier phase of this project, most of the required physics (general
relativity, inflationary dynamics, cosmological perturbation theory) was not available
from prior coursework and was learned specifically for this project, primarily from
Baumann's notes, with regular mentor meetings providing direction and correction. A
substantial part of this later phase was methodological: extending a validated
single-potential notebook to eight potentials and six parameter sweeps required
imposing shared conventions (a common $N_{\rm end}$, a common calibration target, a
common tolerance) that the original single-potential version did not need, and several
early attempts at cross-potential comparison surfaced inconsistencies (such as the
$N_{\rm end}=60$ vs.\ $65$ mismatch between the single-potential and sweep notebooks)
that had to be tracked down and resolved before the comparison was meaningful. Working
through the classification hypothesis in Sec.~\ref{sec:classification} was also a
lesson in distinguishing a genuine numerical pattern from an over-eager conclusion: an
earlier, simpler version of the hypothesis was disproved by the non-monotonic $r(g)$
behaviour of the quadratic+quartic mixture, which is why the classification is presented
here explicitly as preliminary rather than as a finished result.

\section{Conclusion}
This project extended a validated single-field numerical inflation pipeline from one
potential (quadratic, scalar-only) to eight potentials with both scalar and tensor
perturbations, all calibrated to a shared $N_{\rm end}=65$, $\mathcal{P}_\mathcal{R}(k_*)=
2.05\times10^{-9}$ benchmark. The scalar sector remains validated against closed-form
quadratic slow-roll expectations to sub-percent accuracy; the tensor sector and the
single-field consistency relation $r=-8n_t$ were checked across all eight potentials and
show violations that track a potential's $|\eta_v/\varepsilon_v|$ at the pivot, largest
for the hilltop-type potentials. Systematic parameter sweeps reveal a classification of
a potential's free parameters into scale-type and shape-type, with distinct signatures in
the $(n_s,r)$ plane, observed and numerically verified across all seven sweeps -- this is
the leading candidate for the project's next phase, with an analytic derivation of the
pattern as the natural next step.

\begin{thebibliography}{9}
\bibitem{baumann} D.~Baumann, \emph{Cosmology}, lecture notes. Primary theoretical
reference throughout, particularly the sections on inflation, cosmological perturbation
theory, and quantum initial conditions.
\end{thebibliography}

\end{document}
